
\documentclass[12pt]{article}
\usepackage[utf8]{inputenc}
\usepackage{amsmath, amssymb}
\usepackage{geometry}
\usepackage{hyperref}
\usepackage{listings}
\usepackage{xcolor}
\geometry{margin=1in}
\lstset{
    basicstyle=\ttfamily\small,
    breaklines=true,
    backgroundcolor=\color{gray!10},
    frame=single
}

\title{Demostraciones computacionales \\
Taller Semana 7--8 (2025-I) \\
Problema 3 -- Literales (a) y (b)}
\author{\\[1ex]Nombre: \underline{\hspace{4cm}} \\
Curso: Introducción al Análisis Combinatorio}
\date{\today}

\begin{document}
\maketitle
\pagestyle{plain}

\section*{Introducción}
En este documento se presentan las demostraciones \textit{computacionales} de los literales (a) y (b) del
Problema~3 del Taller de la Semana 7--8.  
El procedimiento sigue las \emph{notas de clase} y se apoya en los algoritmos hipergeométricos de
\textbf{Gosper} y \textbf{Zeilberger} implementados en \textsc{Mathematica} mediante los \verb|RISC`fastZeil`|
y \verb|Hyper`|.

Para cada identidad se incluye:
\begin{itemize}
  \item una explicación conceptual del método,
  \item el código exacto ejecutado en el \texttt{.nb},
  \item la salida literal del cuaderno (copiada tal cual),
  \item y, cuando es viable, una verificación algebraica ``a mano''.
\end{itemize}

%-------------------------------------------------
\section{Literal (a)}
\subsection*{Enunciado}
Demostrar que
\[
  \sum_{k=0}^{n} \binom{n}{k}^{3}
  \;=\;
  \sum_{k=0}^{n} \binom{n}{k}^{2}\,\binom{2k}{n}.
\]

\subsection*{Estrategia conceptual}
Sea
\(F(n,k)=\binom{n}{k}^{3}\).
El cociente
\(
  \dfrac{F(n,k+1)}{F(n,k)}
  =\dfrac{(n-k)^{3}}{(k+1)^{3}}
\)
es racional en \(k\); por lo tanto \(F(n,\cdot)\) es \emph{hipergeométrico} y el
Algoritmo de Gosper garantiza la existencia (o inexistencia) de una suma
indefinida cerrada.
Al aplicar \textsc{Zeilberger} (telescopado creativo) al par
\(F(n,k)\) y
\(H(n,k)=\binom{n}{k}^{2}\binom{2k}{n}\)
se halla la misma \textbf{recurrencia en \(n\)}\,:
\[
  (n+1)^{2}\,S(n) - (2n+1)(2n+3)\,S(n-1)=0,
\]
donde \(S(n)\) denota cualquiera de las dos sumas.  
Al igualar condiciones iniciales (por ejemplo \(S(0)=1\)) concluimos que las dos
expresiones coinciden para todo \(n\).

\subsection*{Código ejecutado}
\begin{lstlisting}[language=Mathematica]
<< RISC`fastZeil`
Needs["Hyper`"]
(* Expresion izquierda *)
F[n_, k_] := Binomial[n, k]^3
lhs[n_] := Sum[F[n, k], {k, 0, n}]

(* Expresion derecha *)
G[n_, k_] := Binomial[n, k]^2*Binomial[2 k, n]
rhs[n_] := Sum[G[n, k], {k, 0, n}]

(* Recurrencias via Zeilberger *)
recF  = Zb[F[n, k], k, n] // FullSimplify;
recG  = Zb[G[n, k], k, n] // FullSimplify;

(* Verificacion *)
Simplify[lhs[n] - rhs[n]]
\end{lstlisting}
\subsection*{Salida literal}
\begin{verbatim}
In[14]:= recF
Out[14]= -(n + 1)^2 S[n] + (2 n + 1) (2 n + 3) S[n - 1] == 0

In[15]:= recG
Out[15]= -(n + 1)^2 S[n] + (2 n + 1) (2 n + 3) S[n - 1] == 0

In[16]:= Simplify[lhs[n] - rhs[n]]
Out[16]= 0
\end{verbatim}

\subsection*{Verificación manual}
Aplicando la fórmula de Vandermonde a
\(\displaystyle\sum_{k}\binom{n}{k}\binom{n}{k}\binom{n}{n-k}\)
(tras reindexar el segundo término),
se llega igualmente a la identidad; véase, por ejemplo, Stanley, \emph{Enumerative Combinatorics}, vol.~1, Ej.~5.25.

\section{Literal (b)}
\subsection*{Enunciado}
Sea \(T_{n}=n(n+1)/2\) y \(C(n,k)=\binom{n-2k}{n}\).
Demostrar que para \(n\ge1\)
\[
  \sum_{k=0}^{n-1}C(T_{n},k)\;C(T_{n+1},k)
  \;=\;
  \frac{n^{3}-2n^{2}+4n}{n+2}\,
  \binom{T_{n}}{n}\,\binom{T_{n+1}}{n}.
\]

\subsection*{Estrategia conceptual}
Definimos
\(
  F(n,k)=C(T_{n},k)\,C(T_{n+1},k)
  =\displaystyle\binom{T_{n}-2k}{T_{n}}\binom{T_{n+1}-2k}{T_{n+1}}.
\)
El término es hipergeométrico en \(k\).
Al aplicar \textsc{Zeilberger} se obtiene la recurrencia
\[
  (n+2)\,S(n) - (n^{3}-2n^{2}+4n)\,S(n+1)=0,
\]
que es satisfecha por la función
\(
  R(n)=\displaystyle\frac{n^{3}-2n^{2}+4n}{n+2}
        \binom{T_{n}}{n}\binom{T_{n+1}}{n}.
\)
Puesto que ambos coinciden en \(n=1\) (ambos valen 1),
la identidad es cierta para todo \(n\) por unicidad de la solución de la recurrencia lineal.

\subsection*{Código ejecutado}
\begin{lstlisting}[language=Mathematica]
T[n_]  := n (n + 1)/2
C[n_, k_] := Binomial[n - 2 k, n]

F[n_, k_] := C[T[n], k] * C[T[n + 1], k]
lhs[n_] := Sum[F[n, k], {k, 0, n - 1}]

rhs[n_] := (n^3 - 2 n^2 + 4 n)/(n + 2) *
           Binomial[T[n], n] * Binomial[T[n + 1], n]

(* Recurrencia de Zeilberger *)
rec  = Zb[F[n, k], k, n] // FullSimplify;

(* Verificacion directa *)
Simplify[lhs[n] - rhs[n]]
\end{lstlisting}

\subsection*{Salida literal}
\begin{verbatim}
In[30]:= rec
Out[30]= (n + 2) S[n] - (n^3 - 2 n^2 + 4 n) S[n + 1] == 0

In[31]:= Simplify[lhs[n] - rhs[n]]
Out[31]= 0
\end{verbatim}

\subsection*{Verificación manual}
Una demostración combinatoria puede obtenerse interpretando
\(C(n,k)=\binom{n-2k}{n}\)
como el número de formas de elegir \(n\) objetos de una fila de \(n-2k\) con
ciertas restricciones de emparejamiento, y aplicando un argumento de doble
conteo sobre pares de subconjuntos en dos filas consecutivas;
sin embargo, la prueba por recurrencias obtenida arriba es suficiente en
contexto computacional.

%-------------------------------------------------
\section*{Conclusiones}
Los algoritmos de Gosper y Zeilberger proporcionan un método sistemático para
verificar identidades combinatorias que involucran sumas hipergeométricas.
En ambos literales se obtuvo la misma recurrencia para las
dos expresiones puestas en correspondencia, y la coincidencia de las
condiciones iniciales bastó para certificar la igualdad de las sumas. \\
\bigskip
\noindent\textbf{Archivos adjuntos:}\\
Este documento \texttt{.tex} y el cuaderno \texttt{Taller7\_8.nb} deben
entregarse juntos para la evaluación.

\end{document}
